% Periodic Pattern Learning in the Surrogate Neural Network
% Explains why the CPG's time-varying curvature creates a frequency-domain
% challenge for the surrogate, reviews approaches, and presents solutions.

\documentclass[11pt]{article}

% ---------------------------------------------------------------------------
% Font and encoding
% ---------------------------------------------------------------------------
\usepackage{palatino}
\usepackage[T1]{fontenc}
\usepackage[utf8]{inputenc}

% ---------------------------------------------------------------------------
% Page layout
% ---------------------------------------------------------------------------
\usepackage[top=1in, bottom=1in, left=1.25in, right=1.25in]{geometry}
\usepackage{setspace}
\setstretch{1.15}
\setlength{\parindent}{1em}
\setlength{\parskip}{0.4em}

% ---------------------------------------------------------------------------
% Math
% ---------------------------------------------------------------------------
\usepackage{amsmath}
\usepackage{amssymb}
\usepackage{mathtools}
\usepackage{bm}

% ---------------------------------------------------------------------------
% Figures and tables
% ---------------------------------------------------------------------------
\usepackage{booktabs}
\usepackage{float}
\usepackage{array}
\usepackage{tabularx}
\usepackage{multirow}

% ---------------------------------------------------------------------------
% Colors and hyperlinks
% ---------------------------------------------------------------------------
\usepackage{xcolor}
\usepackage[hidelinks]{hyperref}

% ---------------------------------------------------------------------------
% Lists
% ---------------------------------------------------------------------------
\usepackage{enumitem}
\setlist[itemize]{nosep, topsep=0.3em, itemsep=0.2em}

% ---------------------------------------------------------------------------
% Section formatting
% ---------------------------------------------------------------------------
\usepackage{titlesec}
\titlespacing*{\section}{0pt}{1.2em}{0.6em}
\titlespacing*{\subsection}{0pt}{0.9em}{0.3em}

% ---------------------------------------------------------------------------
% Custom commands (consistent with system-formulation.tex)
% ---------------------------------------------------------------------------
\newcommand{\vect}[1]{\mathbf{#1}}
\newcommand{\mat}[1]{\mathbf{#1}}
\newcommand{\Real}{\mathbb{R}}

% ---------------------------------------------------------------------------
\begin{document}
% ---------------------------------------------------------------------------

\begin{center}
  {\Large\bfseries Periodic Pattern Learning in the Surrogate Network} \\[0.5em]
  {\large Spectral Bias, Fourier Features, and Input Encoding for the CPG-Driven Cosserat Rod} \\[1.5em]
  \textit{Snake-HRL Project --- Reference Document}
\end{center}

\vspace{1em}

% ===================================================================
\section{The Problem}
\label{sec:problem}
% ===================================================================

The neural surrogate $f_\theta$ replaces the transition operator
$T\colon \Real^{129} \to \Real^{124}$ that composes 500 Position Verlet substeps of a
Cosserat rod simulation.
During those substeps, the CPG (Central Pattern Generator) produces a
\emph{time-varying} rest curvature at each Voronoi node:

\begin{equation}
  \kappa_v^{\mathrm{rest}}(t^n)
  = A \sin\!\bigl(2\pi k \cdot s_v + 2\pi f \cdot t^n + \phi_0\bigr) + b,
  \qquad v = 1, \ldots, 19
  \label{eq:cpg-substep}
\end{equation}

\noindent
where $(A, f, k, \phi_0, b) \in \Real^5$ is the RL action, held fixed for the entire
control interval $\Delta t_{\mathrm{ctrl}} = 0.5$\,s, and
$t^n = t_{\mathrm{step}} + n \,\Delta t$ advances at every substep.
Because of the $2\pi f \cdot t^n$ term, the 500 substeps are \emph{not identical}:
each sees a different curvature configuration, and the rod responds dynamically to
the propagating wave.

The surrogate must therefore learn the \emph{integrated effect} of this
oscillatory forcing:

\begin{equation}
  \vect{s}_{t+1} = T(\vect{s}_t, \vect{a}_t)
  = \underbrace{V_{\Delta t}\!\bigl(\,\cdot\,;\;\bm{\kappa}^{\mathrm{rest},\,499}\bigr)
    \;\circ\; \cdots \;\circ\;
    V_{\Delta t}\!\bigl(\,\cdot\,;\;\bm{\kappa}^{\mathrm{rest},\,0}\bigr)}_{%
    500 \text{ substeps, each with distinct } \bm{\kappa}^{\mathrm{rest}}}
  \;\bigl(\vect{s}_t\bigr)
\end{equation}

\noindent
This is not a power iteration $V_{\Delta t}^{500}$.
The qualitative behaviour of $T$ depends strongly on how many oscillation cycles
fit within $\Delta t_{\mathrm{ctrl}}$:

\begin{equation}
  n_{\mathrm{cycles}} = f \cdot \Delta t_{\mathrm{ctrl}}
  \label{eq:ncycles}
\end{equation}

\begin{table}[H]
  \centering
  \renewcommand{\arraystretch}{1.2}
  \begin{tabular}{ccl}
    \toprule
    $f$ (Hz) & $n_{\mathrm{cycles}}$ & Dynamics character \\
    \midrule
    0.5 & 0.25 & Slow sweep --- rod barely completes a quarter-wave \\
    1.0 & 0.50 & Half-cycle --- rod swings in one direction \\
    2.0 & 1.00 & Full cycle --- rod nearly returns to its initial configuration \\
    3.0 & 1.50 & One-and-a-half cycles --- rod reverses mid-step \\
    \bottomrule
  \end{tabular}
  \caption{Qualitative effect of the CPG frequency on one RL step.
    At integer $n_{\mathrm{cycles}}$, the angular velocity $\omega_z$ nearly
    returns to its initial value; at half-integer $n_{\mathrm{cycles}}$,
    $\omega_z$ reverses sign.}
  \label{tab:freq-regimes}
\end{table}

\noindent
The mapping $f \mapsto \Delta\omega_z$ is therefore \emph{itself oscillatory}:
the change in angular velocity is approximately zero when $n_{\mathrm{cycles}}$
is an integer and maximal when it is a half-integer.
An MLP with standard activations must approximate this sinusoidal relationship
from raw scalar inputs --- precisely the setting where \emph{spectral bias}
becomes problematic.

% ===================================================================
\section{Spectral Bias}
\label{sec:spectral-bias}
% ===================================================================

\textbf{Spectral bias} (the ``frequency principle'') is the well-documented
tendency of neural networks to learn low-frequency components of a target function
before high-frequency components~\cite{rahaman2019spectral}.
Networks with standard activations (ReLU, SiLU, tanh) preferentially fit smooth,
slowly-varying functions.
High-frequency patterns in the target require more training time, more parameters,
or may never be fully learned.

The root cause is that the Neural Tangent Kernel (NTK) of standard MLPs has a rapid
spectral falloff: the kernel's eigenvalues decay with increasing frequency, so
gradient descent converges slowly on high-frequency modes.

\paragraph{Relevance to the surrogate.}
The surrogate's target function contains an oscillatory dependence on the frequency
parameter $f$.
Concretely, $\Delta\omega_z$ as a function of $f$ is approximately sinusoidal
(Table~\ref{tab:freq-regimes}).
An MLP receiving $f$ as a raw normalized scalar in $[-1, 1]$ must learn this
periodicity purely from data.
This is consistent with the poor angular velocity prediction observed empirically
($R^2 = 0.23$ on $\omega_z$ components).

\paragraph{Severity assessment.}
The spectral bias risk is \textbf{localised} to the frequency and wavenumber
dimensions of the action space.
The spatial CPG phase is already handled by the per-element encoding (Section~\ref{sec:existing}),
which provides explicit sinusoidal features.
The remaining gap is the \emph{frequency-domain} relationship between $(f, k)$ and
the integrated dynamical response.

% ===================================================================
\section{The Existing Per-Element Phase Encoding}
\label{sec:existing}
% ===================================================================

The current input to the surrogate is a 193-dimensional vector:

\begin{equation}
  \vect{x} = \bigl[\;
    \underbrace{\vect{s}_{\mathrm{rel}}}_{\Real^{128}}
    \;\big|\;
    \underbrace{\vect{a}}_{\Real^{5}}
    \;\big|\;
    \underbrace{\vect{\gamma}^{\mathrm{phase}}}_{\Real^{60}}
  \;\bigr]
  \;\in \Real^{193}
  \label{eq:current-input}
\end{equation}

\noindent
where $\vect{s}_{\mathrm{rel}}$ is the relative rod state (CoM-subtracted positions,
heading as $\sin/\cos$, velocities, yaw, $\omega_z$), $\vect{a}$ is the raw action,
and $\vect{\gamma}^{\mathrm{phase}}$ contains 60 per-element CPG phase features:

\begin{equation}
  \vect{\gamma}^{\mathrm{phase}} =
  \bigl[\;
    \underbrace{\sin(\psi_j)}_{j=1}^{20}
    \;\big|\;
    \underbrace{\cos(\psi_j)}_{j=1}^{20}
    \;\big|\;
    \underbrace{A \sin(\psi_j)}_{j=1}^{20}
  \;\bigr]
  \;\in \Real^{60}
  \label{eq:phase-enc}
\end{equation}

\noindent
with $\psi_j = 2\pi k \cdot s_j + 2\pi f \cdot t_{\mathrm{step}} + \phi_0$
being the CPG phase angle at element~$j$ at the \emph{start} of the RL step.

\subsection{Strengths}

\begin{enumerate}[label=(\alph*)]
  \item \textbf{Spatially resolved.}
    Each element gets its own phase, capturing the traveling wave's spatial structure.
    The network can distinguish head-to-tail curvature patterns rather than seeing a
    single global phase.

  \item \textbf{$(\sin, \cos)$ pair avoids phase wrapping.}
    The two-component representation is continuous everywhere, avoiding the $2\pi$
    discontinuity that would arise from providing the raw angle~$\psi_j$.

  \item \textbf{Commanded curvature is explicit.}
    The $A \sin(\psi_j)$ features directly provide the rest curvature
    $\kappa_j^{\mathrm{rest}}$ at step start --- the actual forcing term that enters
    the constitutive relation
    $\bm{\tau}_v^L = \mat{B}_{\mathrm{blk}}\,(\bm{\kappa}_v - \bm{\kappa}_v^{\mathrm{rest}})$.

  \item \textbf{Structurally equivalent to Fourier features.}
    The encoding is a hand-crafted Fourier feature expansion of the CPG phase at
    each spatial location, with exactly the right frequency (the CPG fundamental).
    This directly addresses spectral bias for the spatial-temporal phase.
\end{enumerate}

\subsection{Limitations}
\label{sec:existing-limitations}

\begin{enumerate}[label=(\alph*)]
  \item \textbf{Fundamental frequency only.}
    The encoding provides $\sin(\psi_j)$ and $\cos(\psi_j)$ at the CPG fundamental.
    The \emph{integrated response} of 500 nonlinear substeps generates harmonics ---
    the rod's actual curvature at step end is not a pure sinusoid even when the
    forcing is.
    Cross-frequency terms $\sin(2\psi_j)$ and $\sin(\psi_j)\sin(\psi_k)$ are absent.

  \item \textbf{No explicit $n_{\mathrm{cycles}}$ representation.}
    The number of oscillation cycles per step (Eq.~\ref{eq:ncycles}) is the
    critical quantity that determines whether $\omega_z$ increases, decreases, or
    reverses.
    It is available implicitly from $f$ in the action, but not encoded as a periodic feature.

  \item \textbf{Phase at step start only.}
    The encoding provides the CPG state at $t_{\mathrm{step}}$, but the integrated
    response depends on the phase \emph{trajectory} over 500 substeps.
    The phase sweeps from $\psi_j(0)$ to $\psi_j(0) + 2\pi f \cdot \Delta t_{\mathrm{ctrl}}$.
    This sweep information must be inferred from the raw scalar~$f$.
\end{enumerate}

\paragraph{Verdict.}
The existing encoding is well-designed but has a specific gap: the network must
learn the oscillatory relationship between $f$ and the integrated dynamical response
from a raw scalar.
Fourier-encoding $f$ (and $k$) directly targets this gap.

% ===================================================================
\section{Approaches Considered and Not Recommended}
\label{sec:not-recommended}
% ===================================================================

\subsection{SIREN (Periodic Activation Functions)}

SIREN~\cite{sitzmann2020siren} replaces all activation functions with
$\sin(\omega_0 \cdot x)$, where $\omega_0$ is a frequency hyperparameter
(typically 30).
Its derivatives are also sinusoidal, making it effective for learning signals
whose derivatives must be accurate (e.g., signed distance functions, PDE solutions).

\paragraph{Why not here.}
The periodicity in the surrogate's problem is in the \emph{input encoding},
not in the \emph{target function}.
The mapping from (body shape + CPG phase features) to (body shape delta)
is not itself periodic --- it is a complex nonlinear regression where the periodic
structure has already been factored out via the phase encoding.
SIREN also produces ringing artifacts when the target contains non-periodic
components (position drift, friction effects), and requires a specialised
initialisation scheme that may interact poorly with the existing LayerNorm + residual
block architecture.

\subsection{Snake Activation}

The Snake activation~\cite{ziyin2020neural}
$\mathrm{Snake}_a(x) = x + \tfrac{\sin^2(ax)}{a}$ combines ReLU-like monotonicity
with a periodic inductive bias.
Its primary advantage is \emph{extrapolation}: given periodic training data,
Snake networks extend the pattern beyond the training domain.

\paragraph{Why not here.}
The surrogate operates within a bounded action space and state distribution.
Extrapolation is not the primary concern --- the body-frame representation and
dense training coverage address generalisation more fundamentally.
Switching activations introduces risk (different initialisation, learning rate
sensitivity) with minimal expected gain since the periodicity is already in the input.

\subsection{Neural ODEs}

Neural ODEs~\cite{chen2018neuralode} learn continuous-time dynamics
$d\vect{x}/dt = f_\theta(\vect{x}, t)$ and integrate via an ODE solver.
They are appropriate when the time horizon varies or adaptive step sizes are needed.

\paragraph{Why not here.}
The surrogate is a \emph{fixed-step, fixed-horizon} map
($\Real^{193} \to \Real^{128}$).
The entire point is to replace 500 integration steps with a single forward pass.
A Neural ODE would reintroduce stepwise integration in the forward pass,
defeating the purpose.

\subsection{Fourier Neural Operators (FNO)}

FNO~\cite{li2021fno} learns mappings between function spaces via spectral
convolutions on spatial grids, excelling at parametric PDEs with resolution invariance.

\paragraph{Why not here.}
The surrogate maps between \emph{finite-dimensional vectors}, not functions.
The 20-element rod is already discretised at a fixed resolution.
FNO is architecturally mismatched --- it would add complexity without benefit for
this specific problem.

% ===================================================================
\section{Recommended Solutions}
\label{sec:solutions}
% ===================================================================

The recommendations target the specific gap identified in
Section~\ref{sec:existing-limitations}: the raw scalar representation of the
frequency~$f$ and wavenumber~$k$ in the action.

% -------------------------------------------------------------------
\subsection{Solution 1: Random Fourier Features on $(f, k)$}
\label{sec:rff}
% -------------------------------------------------------------------

\paragraph{Theory.}
Tancik et al.~\cite{tancik2020fourier} showed that mapping low-dimensional inputs
through sinusoidal functions before feeding them to an MLP eliminates spectral bias.
The encoding is:

\begin{equation}
  \gamma(\vect{v}) =
  \bigl[\;
    \sin\!\bigl(2\pi \mat{B}\,\vect{v}\bigr),\;
    \cos\!\bigl(2\pi \mat{B}\,\vect{v}\bigr)
  \;\bigr]
  \;\in \Real^{2m}
  \label{eq:rff}
\end{equation}

\noindent
where $\vect{v} \in \Real^d$ is the input,
$\mat{B} \in \Real^{m \times d}$ is a random frequency matrix with
$B_{ij} \sim \mathcal{N}(0, \sigma^2)$, and
$\sigma$ controls the bandwidth (the range of frequencies represented).
Each row of $\mat{B}$ defines a random frequency direction in input space;
the $(\sin, \cos)$ pair at that frequency provides the network with a basis
function at the corresponding scale.

The connection to the NTK is direct: Fourier features reshape the NTK's spectrum
to have uniform eigenvalues across a wider frequency band, removing the
low-frequency preference.

\paragraph{Application.}
Apply Eq.~\ref{eq:rff} to the frequency and wavenumber components of the action:

\begin{equation}
  \vect{v} = \bigl(f_{\mathrm{norm}},\; k_{\mathrm{norm}}\bigr) \in [-1, 1]^2
  \qquad
  \gamma(\vect{v}) \in \Real^{2m}
  \label{eq:rff-fk}
\end{equation}

\noindent
with $m = 32$ (yielding 64 additional features) and $\mat{B} \in \Real^{32 \times 2}$.
The bandwidth $\sigma$ is a hyperparameter to be swept.

\paragraph{Bandwidth selection.}
The target function's characteristic frequency in normalised input space determines
the optimal~$\sigma$.
The mapping $f \mapsto \Delta\omega_z$ oscillates on the scale of
$1 / \Delta t_{\mathrm{ctrl}} = 2$\,Hz in physical units, corresponding to a
normalised frequency of $\mathcal{O}(1\text{--}10)$ in $[-1, 1]$ space.
A sweep over $\sigma \in \{1, 3, 10, 30\}$ covers the plausible range:

\begin{table}[H]
  \centering
  \renewcommand{\arraystretch}{1.2}
  \begin{tabular}{cl}
    \toprule
    $\sigma$ & Effect \\
    \midrule
    1   & Low-frequency features only --- may underfit the oscillatory dependence \\
    3   & Moderate --- captures the primary oscillation \\
    10  & High-frequency --- resolves fine structure in $f \mapsto \Delta\omega_z$ \\
    30  & Very high-frequency --- risk of overfitting to noise \\
    \bottomrule
  \end{tabular}
  \caption{Fourier feature bandwidth $\sigma$ sweep.
    Expected optimum: $\sigma \in [3, 10]$.}
  \label{tab:sigma-sweep}
\end{table}

\paragraph{Implementation.}
A \texttt{FourierFeatures} module samples $\mat{B}$ once at initialisation
and registers it as a non-trainable buffer (or optionally as a learnable parameter):

\begin{equation}
  \boxed{\;
  \begin{aligned}
    \mat{B} &\sim \mathcal{N}(\vect{0},\; \sigma^2 \mat{I}),
      \quad \mat{B} \in \Real^{m \times 2}
      && [\text{sample once at init}] \\[4pt]
    \vect{p} &= 2\pi\, \mat{B}\, \vect{v}
      && [\text{projection, } \vect{p} \in \Real^m] \\[4pt]
    \gamma(\vect{v}) &= \bigl[\sin(\vect{p}),\; \cos(\vect{p})\bigr]
      && [\text{output, } \gamma \in \Real^{2m}]
  \end{aligned}
  \;}
  \label{eq:rff-impl}
\end{equation}

\noindent
The extended surrogate input becomes:

\begin{equation}
  \vect{x}' = \bigl[\;
    \vect{s}_{\mathrm{rel}}
    \;\big|\;
    \vect{a}
    \;\big|\;
    \vect{\gamma}^{\mathrm{phase}}
    \;\big|\;
    \gamma(f_{\mathrm{norm}}, k_{\mathrm{norm}})
  \;\bigr]
  \;\in \Real^{193 + 2m}
  \label{eq:extended-input}
\end{equation}

\noindent
For $m = 32$: $\dim(\vect{x}') = 257$.

\paragraph{Why this is the highest-impact change.}
The Fourier features directly provide the sinusoidal basis functions that the
MLP would otherwise need to synthesise from the raw scalar $f$.
Instead of learning $\Delta\omega_z \approx g(\sin(c \cdot f))$ for some
unknown constant $c$, the network receives $\sin(2\pi b_i \cdot f)$ at
multiple frequencies $b_i$ and need only learn a \emph{linear combination}
of these basis functions.
This converts a hard nonlinear approximation problem into a substantially
easier regression.

% -------------------------------------------------------------------
\subsection{Solution 2: Explicit $n_{\mathrm{cycles}}$ Feature}
\label{sec:ncycles}
% -------------------------------------------------------------------

Add a derived feature that directly encodes how many oscillation cycles the CPG
completes within one RL step:

\begin{equation}
  n_{\mathrm{cycles}} = f \cdot \Delta t_{\mathrm{ctrl}}
  \qquad
  \gamma_{\mathrm{cyc}} = \bigl[\sin(2\pi\, n_{\mathrm{cycles}}),\;
                                 \cos(2\pi\, n_{\mathrm{cycles}})\bigr]
  \;\in \Real^2
  \label{eq:ncycles-enc}
\end{equation}

\noindent
where $f$ is the denormalised frequency in Hz.
When $n_{\mathrm{cycles}}$ is an integer (full cycles), $\sin(2\pi n_{\mathrm{cycles}}) = 0$
and $\cos(2\pi n_{\mathrm{cycles}}) = 1$; when it is a half-integer (reversal),
$\sin = 0$ and $\cos = -1$.
This directly tells the network whether the step ends mid-oscillation.

The cost is 2 additional input dimensions.
The frequency must be denormalised from the action before computing $n_{\mathrm{cycles}}$:

\begin{equation}
  f = f_{\min} + \frac{f_{\mathrm{norm}} + 1}{2}\,(f_{\max} - f_{\min})
  = 0.5 + \frac{f_{\mathrm{norm}} + 1}{2} \cdot 2.5
  \quad [\text{Hz}]
\end{equation}

This is the simplest change with a clear physical motivation.

% -------------------------------------------------------------------
\subsection{Solution 3: Harmonic Phase Features}
\label{sec:harmonics}
% -------------------------------------------------------------------

The nonlinear rod dynamics generate harmonics of the CPG fundamental.
Even though the forcing is sinusoidal at frequency $k$, the rod's curvature
response contains $2k$, $3k$, and higher harmonics due to:

\begin{itemize}
  \item Geometric nonlinearity (finite rotations couple curvature modes)
  \item Quadratic velocity terms in the damping forces
  \item Contact and friction interactions
\end{itemize}

\noindent
The current encoding provides only the fundamental.
Extending to $H$ harmonics:

\begin{equation}
  \vect{\gamma}^{\mathrm{harm}} =
  \bigl[\;
    \underbrace{\sin(h\,\psi_j),\; \cos(h\,\psi_j)}_{\text{for each } h = 1, \ldots, H}
    \;\big|\;
    A \sin(\psi_j)
  \;\bigr]_{j=1}^{20}
  \;\in \Real^{20(2H + 1)}
  \label{eq:harmonics}
\end{equation}

\noindent
For $H = 3$ (fundamental + 2nd + 3rd harmonic):
$\dim(\vect{\gamma}^{\mathrm{harm}}) = 20 \times 7 = 140$,
adding 80 features beyond the current 60.

The $h$-th harmonic provides a basis function that oscillates $h$ times as fast
along the rod body and in time, letting the network represent the nonlinear
response without synthesising these harmonics internally.

% ===================================================================
\section{Combined Input Architecture}
\label{sec:combined}
% ===================================================================

Combining all three solutions, the enhanced surrogate input is:

\begin{equation}
  \vect{x}^{\star} = \bigl[\;
    \underbrace{\vect{s}_{\mathrm{rel}}}_{\Real^{128}}
    \;\big|\;
    \underbrace{\vect{a}}_{\Real^{5}}
    \;\big|\;
    \underbrace{\vect{\gamma}^{\mathrm{harm}}}_{\Real^{140}}
    \;\big|\;
    \underbrace{\gamma(f, k)}_{\Real^{64}}
    \;\big|\;
    \underbrace{\gamma_{\mathrm{cyc}}}_{\Real^{2}}
  \;\bigr]
  \;\in \Real^{339}
  \label{eq:full-input}
\end{equation}

\noindent
Compared to the current 193-dimensional input, this adds 146 features.
The MLP hidden dimension (512) and depth (3 layers) remain unchanged;
only the input projection layer widens.

\begin{table}[H]
  \centering
  \renewcommand{\arraystretch}{1.3}
  \begin{tabular}{llcc}
    \toprule
    \textbf{Component} & \textbf{What it encodes} & \textbf{Current} & \textbf{Proposed} \\
    \midrule
    Relative state $\vect{s}_{\mathrm{rel}}$
      & Rod shape, velocities, heading & 128 & 128 \\
    Action $\vect{a}$
      & Raw CPG parameters (normalised) & 5 & 5 \\
    Per-element phase
      & CPG phase at each element (fundamental) & 60 & --- \\
    Harmonic phase $\vect{\gamma}^{\mathrm{harm}}$
      & CPG phase at each element ($H = 3$ harmonics) & --- & 140 \\
    Fourier features $\gamma(f, k)$
      & Frequency-domain encoding of $f$ and $k$ & --- & 64 \\
    Cycle encoding $\gamma_{\mathrm{cyc}}$
      & $\sin/\cos$ of oscillation cycles per step & --- & 2 \\
    \midrule
    \textbf{Total} & & \textbf{193} & \textbf{339} \\
    \bottomrule
  \end{tabular}
  \caption{Surrogate input dimensions: current vs.\ proposed.
    The harmonic phase features ($H = 3$) subsume the current fundamental-only
    encoding ($H = 1$).}
  \label{tab:input-dims}
\end{table}

\paragraph{Parameter count impact.}
The first hidden layer grows from $193 \times 512 = 98{,}816$ to
$339 \times 512 = 173{,}568$ parameters (a 76\% increase in the first layer).
Total model parameters increase from $\sim$790k to $\sim$860k ---
a modest 9\% increase.
All subsequent layers are unchanged.

% ===================================================================
\section{Experiment Design}
\label{sec:experiments}
% ===================================================================

The three solutions should be tested \emph{independently} before combining,
to isolate the contribution of each.

\begin{table}[H]
  \centering
  \renewcommand{\arraystretch}{1.3}
  \begin{tabular}{clccc}
    \toprule
    \textbf{Priority} & \textbf{Experiment} & \textbf{Extra dims} & \textbf{Expected gain} & \textbf{Risk} \\
    \midrule
    1 & $n_{\mathrm{cycles}}$ encoding (Eq.~\ref{eq:ncycles-enc})
      & +2  & Medium & None \\
    2 & Fourier features on $(f, k)$, $\sigma$ sweep (Eq.~\ref{eq:rff-fk})
      & +64 & High & Low \\
    3 & Harmonic phase features, $H = 3$ (Eq.~\ref{eq:harmonics})
      & +80 & Medium & Low \\
    4 & Combined: all three
      & +146 & High & Low \\
    \bottomrule
  \end{tabular}
  \caption{Experiment priority.
    Primary metric: validation $R^2$ on $\omega_z$ components
    (current baseline: 0.23).
    Secondary metric: validation MSE on all 124 state dimensions.}
  \label{tab:experiments}
\end{table}

\paragraph{Fourier feature $\sigma$ sweep.}
For experiment~2, sweep $\sigma \in \{1, 3, 10, 30\}$ with $m = 32$ fixed.
All other hyperparameters (learning rate, batch size, epochs) remain at their
current values to isolate the effect of the encoding.

\paragraph{Ablation.}
If the combined model (experiment~4) improves $R^2$, ablate by removing one
component at a time to confirm each contributes independently.

% ===================================================================
\section{Why Not Fourier-Encode Everything?}
\label{sec:why-not-all}
% ===================================================================

A natural question: if Fourier features help for $(f, k)$, why not apply them
to all 128 state dimensions?

\begin{enumerate}[label=(\arabic*)]
  \item \textbf{Curse of dimensionality.}
    Random Fourier features scale as $\mathcal{O}(2m \cdot d)$ in output dimension.
    For $d = 128$ and $m = 32$, this would add $128 \times 64 = 8{,}192$ features ---
    an input dimension explosion that dominates the model.

  \item \textbf{The state is not periodic.}
    Spectral bias is a problem when the \emph{target function} oscillates as a
    function of the \emph{input coordinates}.
    The state vector (positions, velocities, yaw angles) does not enter the target
    as periodic arguments.
    The target's dependence on the state is smooth and well-captured by standard
    activations.

  \item \textbf{The bottleneck is localised.}
    The poor $\omega_z$ prediction is specifically attributable to the
    frequency-dependent oscillatory dynamics (Section~\ref{sec:problem}).
    Fourier features on $(f, k)$ target this bottleneck precisely.
    Applying them elsewhere dilutes the signal.
\end{enumerate}

% ===================================================================
\section{Summary}
\label{sec:summary}
% ===================================================================

The CPG's time-varying rest curvature creates a frequency-domain challenge for
the surrogate: the integrated effect of 500 substeps is an oscillatory function
of the action parameters $(f, k)$, and standard MLPs learn such functions slowly
(spectral bias).

The existing per-element phase encoding already addresses the
\emph{spatial-temporal phase} at each element --- it is a well-designed Fourier
feature expansion of the CPG state.
But it does not address the \emph{frequency-domain} relationship between
$(f, k)$ and the integrated dynamical response.

Three targeted solutions are proposed:

\begin{enumerate}
  \item \textbf{Random Fourier Features on $(f, k)$} (Section~\ref{sec:rff}):
    Provides sinusoidal basis functions of the frequency parameters, converting
    a hard nonlinear approximation problem into an easier regression.
    Highest expected impact.

  \item \textbf{Explicit $n_{\mathrm{cycles}}$ encoding} (Section~\ref{sec:ncycles}):
    Two features that directly encode how many oscillation cycles fit in the
    control interval. Trivial to implement, physically motivated.

  \item \textbf{Harmonic phase features} (Section~\ref{sec:harmonics}):
    Higher harmonics of the per-element phase, capturing the rod's nonlinear
    response to sinusoidal forcing.
\end{enumerate}

\noindent
SIREN, Snake activations, Neural ODEs, and Fourier Neural Operators were
considered and rejected: the periodicity is in the \emph{input encoding}, not
the \emph{network architecture}, and the surrogate's fixed-step,
finite-dimensional structure does not match these approaches' strengths.

% ===================================================================
% References
% ===================================================================

\begin{thebibliography}{9}

\bibitem{rahaman2019spectral}
N.~Rahaman, A.~Baratin, D.~Arpit, F.~Draxler, M.~Lin, F.~Hamprecht,
Y.~Bengio, and A.~Courville,
``On the spectral bias of neural networks,''
in \textit{Proc.\ ICML}, 2019.

\bibitem{tancik2020fourier}
M.~Tancik, P.~Srinivasan, B.~Mildenhall, S.~Fridovich-Keil, N.~Raghavan,
U.~Singhal, R.~Ramamoorthi, J.~Barron, and R.~Ng,
``Fourier features let networks learn high frequency functions in low
dimensional domains,''
in \textit{Proc.\ NeurIPS}, 2020.

\bibitem{sitzmann2020siren}
V.~Sitzmann, J.~N.~P.~Martel, A.~W.~Bergman, D.~B.~Lindell, and G.~Wetzstein,
``Implicit neural representations with periodic activation functions,''
in \textit{Proc.\ NeurIPS}, 2020.

\bibitem{ziyin2020neural}
L.~Ziyin, T.~Hartwig, and M.~Ueda,
``Neural networks fail to learn periodic functions and how to fix it,''
in \textit{Proc.\ NeurIPS}, 2020.

\bibitem{chen2018neuralode}
R.~T.~Q.~Chen, Y.~Rubanova, J.~Bettencourt, and D.~Duveniste,
``Neural ordinary differential equations,''
in \textit{Proc.\ NeurIPS}, 2018.

\bibitem{li2021fno}
Z.~Li, N.~Kovachki, K.~Azizzadenesheli, B.~Liu, K.~Bhattacharya,
A.~Stuart, and A.~Anandkumar,
``Fourier neural operator for parametric partial differential equations,''
in \textit{Proc.\ ICLR}, 2021.

\end{thebibliography}

\end{document}
